\section{The algebra construction}\label{sec:algebra}

We now prove Theorem~\ref{thm:algebra-package}.  The construction has two
parts.  A pair of monoids supplies a simple algebra and an embedding into
a finitely presented algebra.  Partial affine maps then give a corner
with two equivalent orthogonal subcorners.  A separator in the
monoid construction will prove the tensor detection property.

\subsection{The monoid input}

We first define the auxiliary monoid used in the construction.  Let
\(\mathcal P\) be the monoid of partial maps of \(\mathbb Q\) generated by
the full affine maps
\[
 u(x)=x/2,\qquad v(x)=(x+1)/2,
\]
their inverses, and the partial identities \(e_0,e_1\) on
\([0,\infty)\cap\mathbb Q\) and \([1,\infty)\cap\mathbb Q\), respectively.
Composition applies the right-hand factor first.  The identity
\(1_{\mathcal P}\) is the full identity map.

\begin{lemma}\label{lem:partial-affine}
The monoid \(\mathcal P\) has decidable word problem.  Each of its elements
is a map \(x\mapsto ax+b\) on a nonempty rational tail
\([r,\infty)\cap\mathbb Q\), or on all of \(\mathbb Q\), where
\(a=2^m\) for some \(m\in\mathbb Z\), and \(b,r\) are dyadic rationals.
\end{lemma}

\begin{proof}
Represent a full domain by \(r=-\infty\).  The composition of the data
\((a,b,r)\) and \((c,d,s)\) is
\[
 \left(ac,ad+b,\max\{s,(r-d)/c\}\right).
\]
The generators have such data, and the displayed operation preserves
them.  Each domain is infinite, so two represented maps agree precisely
when their domains, slopes and intercepts agree.  The data and their
comparison are computable using dyadic arithmetic.
\end{proof}

For a monoid \(K\) with zero \(\Omega\ne1_K\), write
\[
 T=(K\setminus\{\Omega\})\times\mathcal P\ \sqcup\ \{0_T\}.
\]
The element \(0_T\) is absorbing.  Nonzero pairs multiply
coordinatewise, except that a product whose \(K\)-coordinate is
\(\Omega\) is replaced by \(0_T\).  Its identity is
\((1_K,1_{\mathcal P})\).

\begin{theorem}\label{thm:monoids}
Let \(\widehat G\) be a finitely generated group with decidable word
problem.  There exist a finitely presented monoid \(K\) with zero
\(\Omega\ne1_K\), and a finitely generated monoid \(H\) with decidable
word problem and zero \(\zeta\ne1_H\), with the following properties.
\begin{enumerate}[label=\textup{(\roman*)}]
\item There is an injective group homomorphism
\(\widehat G\longrightarrow H^\times\).
\item For the monoid \(T\) defined above, there is an injective unital
monoid homomorphism \(\kappa:T\longrightarrow H\) with
\(\kappa(0_T)=\zeta\).
\item For every finite ordered tuple \(h_1,\ldots,h_d\) of distinct
nonzero elements of \(H\), with \(d\ge1\), there are \(L,R\in H\)
such that
\[
 Lh_1R=1_H,\qquad Lh_iR=\zeta\quad(2\le i\le d).
\]
\item There is an injective multiplicative map \(c:H\longrightarrow K\)
whose image is disjoint from \(\{1_K,\Omega\}\).  We write
\(c_h=c(h)\); this map is not required to preserve identities or zeros.
\item There is an element \(z\in H\) such that the map
\[
 (\tau,\tau')\longmapsto\kappa(\tau)z\kappa(\tau')
\]
is an injection from \((T\setminus\{0_T\})^2\) into
\(H\setminus\{\zeta\}\).
\end{enumerate}
\end{theorem}

The proof of Theorem~\ref{thm:monoids} occupies
Section~\ref{sec:bootstrap}, using the finite compiler proved in
Appendix~\ref{sec:compiler}. In the construction, \(\zeta\) is the zero
symbol and \(z\) is the class of a distinct separator letter
\(\mathtt z\).

\subsection{A simple algebra inside a finitely presented algebra}

Throughout this section \(k=\mathbb F_2\).
For any monoid \(D\) with zero \(0_D\ne1_D\), its contracted monoid
algebra \(k_0[D]\) has \(k\)-basis \(D\setminus\{0_D\}\).  Basis
elements multiply as in \(D\), with a product equal to \(0_D\)
interpreted as the zero vector.  It is a unital algebra, with identity
the basis element \(1_D\).

To prove Theorem~\ref{thm:algebra-package} for a given group \(L\),
apply Theorem~\ref{thm:monoids} with \(\widehat G=L\) and set
\[
 A_0=k_0[H],\qquad B=k_0[K].
\]
Both algebras are nonzero.  We distinguish the monoid elements
\(\zeta,\Omega\) from algebra zero, and distinguish \(c_{\zeta}\)
from all three: it is a nonzero basis element of \(B\).

The monoid requirements now serve distinct purposes. The selectors
in Theorem~\ref{thm:monoids}(iii) isolate one coefficient of a nonzero
element of \(A_0\), forcing simplicity. The maps in (ii) and (iv)
provide algebra embeddings in opposite directions, but treat zero
differently: \(\kappa\) preserves it, whereas the linear extension
of \(c\) requires a correction by \(c_{\zeta}\). Finally, the
separator in (v) keeps the ordered pairs indexing a tensor basis
distinct, which will supply the tensor detector.

\begin{lemma}\label{lem:algebra-embeddings}
The algebra \(A_0\) is simple, the algebra \(B\) is finitely presented,
and there are injective algebra homomorphisms
\[
 i:A_0\longrightarrow B,
 \qquad
 \rho:B\otimes_k k[\mathcal P]\longrightarrow A_0.
\]
The map \(\rho\) is unital.  The map \(i\) is nonunital and has support identity
\(e=i(1_{A_0})\ne1_B\).
\end{lemma}

\begin{proof}
Let \(a=\sum_{j=1}^d\lambda_jh_j\) be a nonzero element of \(A_0\),
with distinct nonzero support elements \(h_j\) and nonzero coefficients
\(\lambda_j\in k\).  Choose the selectors from
Theorem~\ref{thm:monoids}(iii).  In the contracted algebra they give
\(LaR=\lambda_1 1_{A_0}\).  Hence every nonzero two-sided ideal contains
the identity, proving simplicity.  This argument also covers a support
element equal to \(1_H\).

The linear extension of \(c\) alone would not respect the contraction
of the zero class, since \(c_{\zeta}\ne\Omega\). For \(h\ne\zeta\), define
\[
 i(h)=c_h-c_{\zeta}
\]
and extend linearly.  Multiplicativity of \(c\) gives
\[
 c_hc_{h'}=c_{hh'},\qquad
 c_hc_{\zeta}=c_{\zeta}c_h=c_{\zeta},\qquad
 c_{\zeta}^2=c_{\zeta}.
\]
Consequently
\[
 (c_h-c_{\zeta})(c_{h'}-c_{\zeta})=c_{hh'}-c_{\zeta}.
\]
When \(hh'=\zeta\), this is zero, as required.  The vectors \(c_h\)
are distinct nonzero basis elements of \(B\).  In particular, the
coefficient of \(c_h\), for \(h\ne\zeta\), in \(i(a)\) is exactly
the coefficient of \(h\) in \(a\), proving injectivity.  Its support
identity is
\[
 e=c_{1_H}-c_{\zeta},\qquad e^2=e,
 \qquad ei(a)=i(a)e=i(a).
\]
Neither basis vector in this expression is \(1_K\); therefore
\(e\ne1_B\), and indeed \(1_B\notin i(A_0)\).

A finite presentation for \(B\) is obtained from the finite monoid
presentation for \(K\).  We may adjoin a generator \(\Omega\) and a
relation identifying it with a word representing the zero, so that
\(\Omega\) is a presentation letter.  Take the same generators in a free associative
unital \(k\)-algebra, replace every equation \(u=v\) by \(u-v=0\),
and impose \(\Omega=0\).  To verify this presentation, first omit the
last relation.  The ideal of binomial relations is precisely the span
of all differences of words in the same monoid congruence class.
One inclusion follows because each contextual use of a defining
equation belongs to the ideal, and a chain of such uses telescopes.
The reverse inclusion follows by mapping each word to its congruence
class in the free vector space on \(K\).  The quotient thus has basis
\(K\); imposing \(\Omega=0\) removes exactly its absorbing zero
class.  It gives \(B\).  Adding the relation \(2\cdot1=0\) to the same
generators and relations also presents \(B\) as a unital ring.

Finally, nonzero basis pairs give a unital algebra isomorphism
\[
 B\otimes_k k[\mathcal P]\cong k_0[T].
\]
It sends \(b\otimes p\), for basis elements \(b\ne\Omega\) and
\(p\in\mathcal P\), to \((b,p)\); products that acquire
\(K\)-coordinate \(\Omega\) vanish on both sides. The linear extension
of \(\kappa\) gives the required injective unital map \(\rho\).
\end{proof}

Identify \(A=i(A_0)\subseteq B\).  It is a simple algebra with its own
identity \(e\), and Theorem~\ref{thm:monoids}(i) embeds
\(L\) in \(A^\times\): a unit of \(H\) and its inverse
give inverse basis elements in \(A_0\), which \(i\) sends to units
relative to \(e\).  We next construct an injective nonunital
endomorphism of \(B\) whose image lies in \(A\).

\subsection{Two equivalent subcorners}

For a dyadic rational \(a\), let \(e_a\) denote the partial identity
on \([a,\infty)\cap\mathbb Q\).  These maps belong to \(\mathcal P\):
\(u^{-1}v\) is translation by \(1\), its powers and conjugates by
powers of \(u\) give all dyadic translations, and conjugating \(e_0\)
by translation by \(a\) gives \(e_a\).  In the monoid itself,
\begin{equation}\label{eq:tail-identities}
 e_a e_b=e_{\max\{a,b\}},\qquad
 w e_a w^{-1}=e_{w(a)}
\end{equation}
for every available full increasing affine map \(w\).  Thus these
are equalities of basis elements in the abstract algebra
\(k[\mathcal P]\).

We require a nonzero corner that decomposes into two orthogonal copies
of itself. The identities in the next lemma have the rectangular-inverse
form familiar from Leavitt's module-type construction
\cite[Section~3]{Leavitt1962}, with the corner idempotent as identity.
Here we realize them directly using the partial affine maps above.

\begin{lemma}\label{lem:binary-tail-corner}
There is a nonzero idempotent \(q\in k[\mathcal P]\) and elements
\(s'_0,s'_1,t'_0,t'_1\in qk[\mathcal P]q\) such that
\[
 t'_i s'_j=\delta_{ij}q,\qquad
 s'_0t'_0+s'_1t'_1=q.
\]
\end{lemma}

\begin{proof}
Set \(q=e_0-e_1\).  Its two basis elements and \(1_{\mathcal P}\)
are distinct, so \(q\ne0,1\);
Equation~\eqref{eq:tail-identities} gives \(q^2=q\).  Put
\[
 w_0=u,\quad w_1=v,\qquad
 q_0=w_0qw_0^{-1}=e_0-e_{1/2},\quad
 q_1=w_1qw_1^{-1}=e_{1/2}-e_1.
\]
Using the same tail identities and expanding in the monoid basis gives
\[
 q_0^2=q_0,\quad q_1^2=q_1,\quad q_0+q_1=q,
 \quad q_0q_1=q_1q_0=0.
\]
For example,
\(q_0q_1=e_{1/2}-e_1-e_{1/2}+e_1=0\).
In particular \(qq_j=q_jq=q_j\).

Define \(s'_j=w_jq\) and \(t'_j=qw_j^{-1}\).  Then
\(s'_jt'_j=q_j\) and \(t'_js'_j=q\).  The identities
\(s'_j=q_jw_j\) and \(t'_j=w_j^{-1}q_j\), together with
\(qq_j=q_jq=q_j\), show
\[
 qs'_j=s'_jq=s'_j,\qquad qt'_j=t'_jq=t'_j.
\]
Moreover \(t'_jq_j=t'_j\) and \(q_js'_j=s'_j\).  For \(i\ne j\)
we therefore have
\[
 t'_is'_j=t'_iq_iq_js'_j=0.
\]
Finally \(\sum_j s'_jt'_j=q_0+q_1=q\).  All calculations use
Equation~\eqref{eq:tail-identities} and linear operations in
\(k[\mathcal P]\).
\end{proof}

Define
\[
 j:B\longrightarrow A_0,\qquad j(b)=\rho(b\otimes q),
 \qquad \phi=i\circ j:B\longrightarrow B.
\]
These maps are nonunital algebra homomorphisms.  The map \(j\) is multiplicative because \(q^2=q\), and it is injective
because \(q\ne0\) and tensor products are over the field \(k\).
Hence \(\phi\) is injective and \(\phi(B)\subseteq A\).  Let
\[
 p=\phi(1_B),\qquad
 s_l=i\rho(1_B\otimes s'_l),\qquad
 t_l=i\rho(1_B\otimes t'_l)\quad(l=0,1).
\]
The idempotent \(p\) belongs to \(A\), so \(p\ne1_B\).
Applying the two algebra maps to Lemma~\ref{lem:binary-tail-corner}
gives
\begin{equation}\label{eq:B-corner-pairs}
 s_l,t_l\in pBp,\qquad
 t_ls_j=\delta_{lj}p,\qquad s_0t_0+s_1t_1=p.
\end{equation}
These elements commute with \(\phi(B)\).  Indeed,
\[
 (b\otimes q)(1_B\otimes s'_l)=b\otimes s'_l
   =(1_B\otimes s'_l)(b\otimes q),
\]
and the identical calculation holds for \(t'_l\); now apply
\(i\rho\). Thus the required self-embedding and binary pairs
centralizing \(\phi(B)\) have been constructed. It remains to verify detection
for the tensor action.

\subsection{A separating tensor vector}

As in Section~\ref{sec:affine}, set \(R=B\otimes_k B^{\mathrm{op}}\),
acting on \(B\) by \((a\otimes b^{\mathrm{op}})m=amb\), and define
\[
 \psi=\phi\otimes\phi^{\mathrm{op}}.
\]
We need one vector on which \(\psi(r)\) acts nontrivially for every
\(r\ne0\). The separator gives this simultaneously for all
\(r\), by realizing distinct basis tensors as distinct monoid elements.

Let \(W=\rho(B\otimes_k k[\mathcal P])\subseteq A_0\), with basis
\(\{\kappa(\tau):\tau\in T\setminus\{0_T\}\}\).  Part (v) of
Theorem~\ref{thm:monoids} proves that the linear map
\[
 D:W\otimes_k W^{\mathrm{op}}\longrightarrow A_0,
 \qquad D(x\otimes y^{\mathrm{op}})=xzy,
\]
is injective: it sends distinct basis tensors to distinct nonzero basis
elements of \(A_0\).  This includes tensors with an identity factor.
Since \(j(B)\subseteq W\), the composite
\[
 R\xrightarrow{\ j\otimes j^{\mathrm{op}}\ }
 W\otimes_k W^{\mathrm{op}}
 \xrightarrow{\ D\ } A_0
 \xrightarrow{\ i\ } B
\]
is a sequence of injective linear maps.  For
\(r=\sum_{\nu}a_{\nu}\otimes b_{\nu}^{\mathrm{op}}\),
its value is
\[
 \begin{aligned}
 i\bigl(D((j\otimes j^{\mathrm{op}})(r))\bigr)
   &=\sum_{\nu}i(j(a_{\nu}))\,i(z)\,i(j(b_{\nu}))\\
   &=\psi(r)\cdot i(z).
 \end{aligned}
\]
Thus \(b_*=i(z)\) satisfies
\begin{equation}\label{eq:tensor-detection}
 r\ne0\quad\Longrightarrow\quad\psi(r)\cdot b_*\ne0.
\end{equation}
The assertion concerns the action after \(\psi\); its proof uses the
explicit separator vector.

\begin{proof}[Proof of Theorem~\ref{thm:algebra-package}]
The algebra \(B\) is nonzero and finitely presented by
Lemma~\ref{lem:algebra-embeddings}.  Its subalgebra \(A=i(A_0)\) is
simple with identity \(e\), and contains \(L\) in its unit
group.  The map \(\phi\) is injective, has image in \(A\), and has
proper support \(p\).  Equation~\eqref{eq:B-corner-pairs} supplies
the binary pairs centralizing \(\phi(B)\), which give the tensor identities
\eqref{eq:ring-pairs} as explained in Section~\ref{sec:reduction}. Finally,
Equation~\eqref{eq:tensor-detection} supplies the required vector
\(b_*\).  These are all the asserted data.
\end{proof}
